\documentclass[11pt,a4paper]{article}

\usepackage[margin=2.2cm]{geometry}
\usepackage[T1]{fontenc}
\usepackage{lmodern}
\usepackage{microtype}
\usepackage{booktabs}
\usepackage{longtable}
\usepackage{tabularx}
\usepackage{enumitem}
\usepackage{xcolor}
\usepackage{hyperref}
\usepackage{url}

\hypersetup{
  colorlinks=true,
  linkcolor=blue!50!black,
  urlcolor=blue!60!black
}

\definecolor{donegreen}{RGB}{38,120,72}
\definecolor{progressorange}{RGB}{176,104,20}
\definecolor{pendinggray}{RGB}{90,90,90}

\newcommand{\done}{\textcolor{donegreen}{\textbf{Completed}}}
\newcommand{\progress}{\textcolor{progressorange}{\textbf{In progress}}}
\newcommand{\pending}{\textcolor{pendinggray}{\textbf{Pending}}}
\newcommand{\code}[1]{\texttt{#1}}

\title{Progress Report\\
\large Forecast-Grounded Decision Question Answering with Verifiable Reasoning\\
for Natural Gas Pipeline Transient Operation}
\author{Project Development Report}
\date{11 July 2026}

\begin{document}

\maketitle

\section{Executive Summary}

The current development cycle advances Task~1 of the project proposal: converting existing PipeClaw scenario questions and PipeFormer forecasts into verifiable, constraint-aware teacher traces. The previous commit established the PDF-aligned teacher-trace schema, compacted the exported trace for lightweight-model distillation, removed obsolete answer-generation and fallback paths, and extended PipeFormer tool inputs to accept the complete parsed-task contract.

The present uncommitted changes complete the missing engineering-constraint coverage described in Sections~1.4--1.5 of the proposal. They add duration-aware pressure checks, flow-capacity and ramp evidence, linepack recovery and reserve checks, compressor-speed verification, equipment-regulation checks, abnormality warnings, complete intervention and dispatch-priority handling, and stricter grounding of the teacher answer. The mock PipeFormer fixture has also been extended from 14 to 15 variables by adding \code{C\_001\_v002} as a synthetic compressor-speed channel. Its cache, tokenizer, normalization data, topology artifacts, and tiny smoke-test checkpoint were regenerated to keep the fixture internally consistent.

Direct inference now reports a complete verification result with no unresolved output variables. This work does not claim an improvement in PipeFormer forecast accuracy: retraining was limited to the synthetic smoke-test checkpoint required by the changed mock-data dimension.

\section{Baseline from the Previous Commit}

The previous commit is:

\begin{quote}
\code{5110cb767746c672aa53c83257933470c128ffa2}\\
``fixing teacher\_trace generation so it matches the required format stated in the pdf and make the teacher\_trace suitable for distillation. Also make the Pipeformer forecast able to accept all parsed task parameters''
\end{quote}

That commit provided the baseline for the present work:

\begin{itemize}[leftmargin=*]
  \item aligned the top-level record with the Section~1.7 teacher-trace format;
  \item retained the LLM-generated \code{final\_answer} as the teacher target;
  \item removed duplicate tool-level final answers and obsolete deterministic answer generation;
  \item removed the sample-CSV inference fallback;
  \item introduced the seven JSON files required by Section~1.9;
  \item made parsed task parameters available to the PipeFormer tool; and
  \item compacted the trace and removed unnecessary local-path metadata.
\end{itemize}

\section{Work Completed in the Current Change Set}

\subsection{Engineering Constraint Coverage}

\begin{longtable}{p{0.20\textwidth}p{0.58\textwidth}p{0.14\textwidth}}
\toprule
\textbf{Constraint area} & \textbf{Implemented capability} & \textbf{Status} \\
\midrule
\endfirsthead
\toprule
\textbf{Constraint area} & \textbf{Implemented capability} & \textbf{Status} \\
\midrule
\endhead
Pressure & Minimum and maximum pressure, warning and violation nodes, margins, contiguous warning/violation episodes, total and maximum continuous out-of-limit duration, recovery time, and simultaneous end-user warning count. & \done \\
Flow & Segment change magnitude, ramp events, transmission-capacity excursions and duration, boundary-flow change rate, abnormal segments, and widening supply--demand-gap evidence. & \done \\
Linepack & Decline episodes, change rate, minimum linepack, recovery target, recovery sufficiency and time, recovered-within-horizon state, and short-term peak-shaving reserve. & \done \\
Compressor & Load, compression ratio, rotational speed, power, operating-envelope state, and regulation margin to failure thresholds. & \done \\
Equipment regulation & Valve-opening range, pressure-regulator range, and allowable boundary-control adjustment magnitude. & \done \\
Abnormality warning & Abnormal pressure drop, sudden flow change, equipment-state anomaly, and cross-signal potential-leak indication. Missing supply or demand data is reported as not evaluated rather than treated as zero. & \done \\
Human intervention & \code{no\_intervention}, \code{monitoring\_only}, \code{operator\_attention\_required}, and \code{immediate\_intervention\_required}, with incomplete verification handled explicitly. & \done \\
Dispatch priority & Safety and supply assurance, equipment protection, linepack margin, flow stability, then energy and cost. False linepack-decline escalation from reserve-only warnings has been removed. & \done \\
\bottomrule
\end{longtable}

\subsection{Teacher-Trace Reliability and Compactness}

The generator now exports a compact but auditable representation of the forecast, verification, and engineering evidence. Passing checks retain aggregate evidence, while non-pass findings retain relevant values, thresholds, margins, and timing. The LLM prompt and deterministic quality checks prevent unsupported repeatability claims, prevent \code{not\_evaluated} rules from being described as passing, and restrict opaque variable identifiers to returned numerical evidence unless explicit variable metadata exists.

The Task~1 source record remains compatible with the Section~1.7 schema. For Task~2, the recommended training representation is a separate compact SFT projection rather than serialization of the complete archival record. A future exporter should derive task-tagged examples for condition parsing, tool invocation, constraint judgment, evidence extraction, and final-answer generation.

\subsection{PipeFormer Mock-Fixture Completion}

The compressor rule library requires rotational-speed verification. The earlier 14-variable mock fixture did not contain a speed channel, which caused \code{compressor\_rotational\_speed\_limit} to be \code{not\_evaluated} and set \code{quality\_flag} to \code{needs\_review}. The current change set:

\begin{itemize}[leftmargin=*]
  \item adds \code{C\_001\_v002} to the generated mock data and prediction mask;
  \item updates the active mapping, topology attention, tokenizer, normalization statistics, and cache to 15 variables;
  \item updates the model configuration to 15 input/output variables and 11 equipment variables;
  \item regenerates the tiny \code{checkpoint-16}; and
  \item adds the previously undeclared \code{pyarrow} and \code{matplotlib} dependencies needed by cache building and training.
\end{itemize}

The training-result timestamp was changed to an ASCII-safe format to prevent Windows locale failures. Generated cache and tokenizer logs containing absolute local paths were removed from version control and ignored.

\section{Verification Evidence}

\begin{table}[htbp]
\centering
\begin{tabularx}{\textwidth}{lX}
\toprule
\textbf{Check} & \textbf{Result} \\
\midrule
Backend regression suite & 14 tests passed. \\
Diff validation & \code{git diff --check} passed. \\
Active PipeFormer mapping & 15 variables; \code{C\_001\_v002} is present and forecastable. \\
Checkpoint dimensions & \code{input\_dim=15}, \code{output\_dim=15}, \code{equipment\_dims=11}. \\
Direct PipeFormer inference & \code{quality\_flag=pass}, \code{verification\_complete=true}, no unresolved output state variables, and no not-evaluated rules. \\
Compressor-speed rule & Evaluated successfully; the current synthetic forecast produces a warning rather than an unavailable result. \\
Tiny checkpoint training & 16 steps completed; best validation loss approximately 2.436. \\
Detailed evaluation & Token accuracy approximately 0.977 on the tiny synthetic validation fixture. \\
\bottomrule
\end{tabularx}
\end{table}

\section{Alignment with the Project Proposal}

\begin{longtable}{p{0.17\textwidth}p{0.52\textwidth}p{0.22\textwidth}}
\toprule
\textbf{Proposal item} & \textbf{Current state} & \textbf{Assessment} \\
\midrule
\endfirsthead
\toprule
\textbf{Proposal item} & \textbf{Current state} & \textbf{Assessment} \\
\midrule
\endhead
Section 1.3 & Parsed tasks contain case, operating-condition number, boundary conditions, disturbance definition, horizon, attention targets, requested outputs, and verification categories. & \done \\
Sections 1.4--1.5 & The seven-file rule library now covers the specified engineering categories and evidence fields. & \done \\
Section 1.6 & The backend executes parsing, real checkpoint inference, rule verification, evidence extraction, and grounded LLM answer generation. & \done \\
Section 1.7 & The exported record preserves the required top-level teacher-trace fields. & \done \\
Section 1.8 & Schema, grounding, rule consistency, incomplete-verification handling, and hallucination checks exist; full dataset-level quality reporting remains outstanding. & \progress \\
Section 1.9 & Seven constraint JSON files exist. Train/validation/test JSONL splits, schema JSON, quality-report workbook, and statistics tables have not yet been produced at project scale. & \progress \\
Task 2 & The trace contains supervision for the five proposed subtasks, and compact projection requirements are identified. Student-model dataset export, LoRA/QLoRA training, baselines, and evaluation remain outstanding. & \pending \\
Task 3 & No distilled student model has yet been integrated into PipeClaw, so system comparison, engineering metrics, latency/cost evaluation, and ablation studies remain outstanding. & \pending \\
\bottomrule
\end{longtable}

\section{Known Limitations}

\begin{itemize}[leftmargin=*]
  \item Current engineering thresholds are normalized proxy limits. Their physical provenance must be established from station limits, simulator settings, standards, or expert review; otherwise the paper should include threshold sensitivity analysis.
  \item The verified example is a tiny synthetic smoke test and is not evidence of real-network forecasting performance.
  \item The added compressor-speed signal is synthetic and exists to complete the mock contract. It must not be described as SCADA telemetry.
  \item The saved teacher-trace JSON must be regenerated locally after the checkpoint update. Automated regeneration was not performed in this report because it would transmit project-derived inputs to the configured external teacher-model service.
  \item Dataset-scale diversity, holdout design, manual spot-checking, and statistical reporting required by Sections~1.8--1.9 remain to be completed.
\end{itemize}

\section{Recommended Next Steps}

\begin{enumerate}[leftmargin=*]
  \item Regenerate teacher traces with the complete checkpoint and confirm that \code{quality\_flag=pass} for fully supported scenarios.
  \item Expand Task~1 generation across the selected PipeClaw scenario sources, then split by case/scenario before deriving any multi-task projections to prevent leakage.
  \item Produce \code{teacher\_trace\_train.jsonl}, \code{teacher\_trace\_valid.jsonl}, \code{teacher\_trace\_test.jsonl}, \code{teacher\_trace\_schema.json}, the quality-report workbook, and dataset statistics required by Section~1.9.
  \item Implement the compact, task-tagged SFT exporter for the five Task~2 subtasks and retain a smaller set of end-to-end tool-call traces.
  \item Establish threshold provenance or perform sensitivity analysis for all proxy limits.
  \item Fine-tune 7B/14B student baselines using LoRA/QLoRA and evaluate parsing, tool calls, constraint judgment, evidence consistency, intervention labels, JSON validity, and hallucination rate.
  \item Integrate the selected student model into PipeClaw and execute the Task~3 scenario, holdout, failure, and ablation experiments.
\end{enumerate}

\section{Recommended Commit Message}

\begin{verbatim}
Complete PDF-aligned constraint verification and compressor-speed coverage

- implement missing pressure, flow, linepack, compressor, equipment, and
  abnormality evidence
- enforce complete verification and grounded compact teacher traces
- add and retrain the 15-variable mock compressor-speed checkpoint
- fix risk escalation, leak validation, and Windows training portability
\end{verbatim}

\end{document}
